% ============================================================
%  第 9 章  总结与展望
% ============================================================
\section{总结与展望}
\label{sec:conclusion}

\subsection{工作总结}

本文针对当前 Web 应用防火墙普遍采用的正则匹配范式\zh{以文本形态为判定依据}
这一方法论缺陷，提出并实现了基于抽象语法树结构指纹的语义检测方法，
完成了从理论构造到工程系统的完整工作。主要工作归纳如下。

在方法层面，本文形式化定义了 SQL 结构指纹，并证明其同时满足参数无关性、
形态无关性与结构敏感性三条性质。这三条性质的正交组合构成了本方法的
理论基础：攻击者可自由操纵的文本形态不影响指纹，而攻击者必须改变的
语义结构必然导致指纹失配。据此，编码混淆类绕过在\emph{原理上}失效，
而非依赖对已知手法的枚举。针对现有语法树方案\zh{只能判定是否变化、
无法定位变化位置}的不足，本文提出路径签名多重集差分算法，
以线性复杂度实现结构级定位，并将差分结果映射为具有安全语义的攻击分类。

在系统层面，本文设计了网关与 RASP 探针的双层联动架构。该设计的依据是
两层在\emph{可见性}上的本质差异：网关仅能观测参数文本，判定必然是启发式的；
探针可观测已完成拼接的真实语句，判定是确定性的。两层通过统一追踪标识关联，
完整还原攻击证据链。系统另包含基于文档对象模型的跨站脚本语义净化、
双重速率限制、带时间衰减的来源信誉评估与哈希链式防篡改审计日志。

在验证层面，本文采用三轮测试法分离各层防护的独立贡献，
设计 19 项测试用例覆盖 OWASP Top 10 的七个类别。结果表明：
网关可判定类攻击的实时拦截率为 $12/12$；全部 19 项漏洞的代码级修复
验证有效；18 项高危漏洞全部获得防护；检测引擎单元测试通过率 $63/63$；
平均检测时延 0.93~ms，优于设计指标。绕过对抗实验显示，
七种文本变形的结构指纹完全一致，均被稳定检出。

\subsection{创新点总结}

\begin{enumerate}
  \item \textbf{结构指纹方法的形式化构造与性质证明。}
        将\zh{语义结构与文本形态正交}这一直觉转化为可证明的数学性质，
        使方法的有效性建立在原理之上而非经验之上。

  \item \textbf{线性复杂度的结构差分定位算法。}
        以路径签名多重集差分替代复杂度过高的树编辑距离，
        在满足安全检测实际需求的前提下将复杂度降至 $O(n)$，
        并首次为语法树类检测方案提供了结构级的定位能力。

  \item \textbf{基于可见性差异的双层架构论证。}
        明确区分两层的判定性质，避免了将确定性要求强加于信息不足的
        网关层，也避免了将拦截时效性要求强加于位置靠后的探针层。

  \item \textbf{检测过程的完整可解释输出。}
        每次拦截均给出结构证据：变异节点的语法路径、语义分类、
        指纹对照与相似度量化，使安全运营人员能够理解判定依据而非
        仅接受一个分值。
\end{enumerate}

\subsection{局限性分析}

诚实记��方法的边界，是评估其适用范围的前提。本文识别出以下五项局限。

\paragraph{对语句可解析性的依赖。}
结构指纹的构造以成功解析为前提。当遇到本文解析器未覆盖的数据库方言时，
系统降级为词法层特征检测，此时结构级的定位能力丧失。
改进方向是引入多解析器投票机制，或基于形式语法自动生成方言解析器。

\paragraph{基线质量对检测精度的制约。}
若学习期内混入攻击流量，其结构将被误记为合法基线，导致同类攻击被放行。
本文以人工审核环节缓解该风险，但审核本身依赖运营人员的判断力。
改进方向是引入离群点检测，对显著偏离多数样本的候选基线自动标记。

\paragraph{二阶注入的溯源缺失。}
攻击载荷先写入存储、在后续查询时才触发的二阶注入，本文可在触发时刻
检出结构异常，但无法溯源至最初的写入请求。完整解决需要跨请求的
数据流污点标记，这将显著增加系统复杂度。

\paragraph{业务逻辑漏洞的防护边界。}
越权、服务端请求伪造等漏洞的判定需要会话状态或后端上下文，
网关层不可见。本文明确将其归属于应用代码层防线，
这一划分是合理的架构决策，但也意味着系统的防护范围存在明确边界。

\paragraph{加密流量的可见性要求。}
网关须作为传输层安全的终止点方能检测流量内容。在端到端加密的部署形态下，
网关层防护失效，此时仅探针层可用。

\subsection{未来改进方向}

基于上述局限，本文提出四个方向的改进设想。

其一，\textbf{跨请求的污点追踪}。将数据的来源标记随其在系统中的流转持续传播，
使二阶注入可被溯源至最初的污染点。技术路径是扩展 RASP 探针的钩取范围，
在数据写入与读取时分别施加标记与检查。

其二，\textbf{基线的结构泛化}。当前基线以精确的结构指纹匹配，
对高度动态的查询（如可变列的报表查询）会产生大量合法指纹。
改进方向是引入结构模板，将可变部分抽象为通配子树，提升基线的表达能力。

其三，\textbf{检测能力的横向扩展}。本文的方法论——以语义结构而非文本形态
作为判定依据——不限于 SQL。同样的思路可应用于表达式语言注入、
模板注入等其他以\zh{改变解释器语义}为目标的攻击类型。

其四，\textbf{与静态分析的结合}。通过静态分析预先推导各查询点可能产生的
语句结构集合，可减少运行时学习的样本需求，并为基线提供独立的交叉验证来源。

\subsection{成员工作总结与心得}

\subsubsection{组长\quad 黄思凯}

两周的开发过程中，有三点体会值得记录。

\paragraph{方法论的选择先于技术的实现。}
项目初期最重要的决策不是选用何种框架，而是确定\zh{以什么作为检测依据}。
一旦确立了\zh{语义结构正交于文本形态}这一判据，后续的算法设计、
架构划分乃至误报控制策略都获得了统一的评判标准。反之，若从
\zh{如何匹配更多攻击特征}出发，则会陷入规则堆砌的老路。
这让我理解到，工程中的\zh{难}往往不在于实现的复杂度，
而在于是否找到了正确的问题表述。

\paragraph{自测是发现设计缺陷的唯一途径。}
项目在设计文档阶段自认为方案已足够严密，但实际运行后仍暴露出多处问题：
探针的异常抑制配置吞掉了拦截决策，导致\zh{能检出却无法阻断}；
参数探测法因构造脱离真实场景而将普通文本误判为攻击；
哈希链因清空操作未同步重置内存状态而永久断裂。这些缺陷无一能通过
纸面推演发现。这一经历让我对\zh{测试驱动}有了具体而非口号式的理解——
测试不是开发完成后的验收环节，而是发现设计缺陷的主要手段。

\paragraph{安全设计需要诚实面对边界。}
在整理测试结果时，我们一度考虑将业务逻辑类漏洞的\zh{网关未拦截}
表述得更为含蓄，以使数据更好看。但最终选择了明确标注防线归属，
并在报告中专门讨论这一架构决策的合理性。安全工程中，
\emph{清楚地知道系统防不住什么，与知道它能防住什么同等重要}。
一份掩盖边界的防护方案，其危害在于会诱导使用者产生不实的安全感，
而这恰是安全事故的常见诱因。

作为组长，我还体会到分工与一致性的平衡之道：把子系统完整地交给
每位成员独立负责，能激发各自的理解与责任心；但判定逻辑这类牵一发
而动全身的核心约定，必须有人端到端地把关。本次项目的顺利推进
得益于这种\zh{模块自治、核心集中}的组织方式。

\subsubsection{组员\quad 罗锦皓}

我负责安全网关与限流模块的实现。最大的收获是对\zh{启发式判定边界}的
理解：网关只能看到参数文本，无法确知它会被拼接到何处，这决定了网关层的
判定在原理上就是概率性的。实现参数检测时，我最初将参数裸拼进模板语句
解析，导致普通文本被误判为攻击；在与组长讨论后改为将参数作为字符串
字面量注入探针，模拟真实的拼接场景，误报随之消除。这让我认识到
检测方法的构造必须与被检测场景的真实语义保持一致，而不能想当然。

\subsubsection{组员\quad 曾铭}

我负责受控靶场与可视化前端。构建漏洞与安全双实现的过程，迫使我从
攻击者的视角逐个理解每处漏洞的成因——为什么闭合一个引号就能绕过
认证，为什么一个 \code{onerror} 属性就能窃取会话。这种理解深度是
单纯阅读教材无法获得的。前端部分，语法树差分渲染的调试让我体会到
\zh{动效服务于信息}这一原则的意义：注入节点的脉冲提示之所以有冲击力，
正因为它指向了判定结论本身，而非单纯的视觉装饰。

\subsubsection{组员\quad 吴靖翔}

我负责控制台服务与安全测试的执行。印象最深的是三轮测试中发现的
哈希链断裂问题：清空事件后新记录的前置哈希仍指向已删除的记录，
完整性校验从此永久失败。定位并修复这个缺陷让我具体理解了
\zh{持久化状态与内存状态必须同步}这类工程细节的重要性。此外，
执行 19 项用例的三轮验证也让我体会到规范化测试记录的价值——
正是完整的每轮结果对照，使各类漏洞的防线归属得以清晰呈现。

通过本次实践，全体成员共同建立了从威胁建模、方案论证到实验验证的
完整方法认知，并对\zh{安全第一、预防为主}这一原则形成了
基于实践的理解。
